\chapter{Introduction}

Software testing is essential for ensuring that changes to a program do not
introduce new faults.  As systems grow in size and complexity, their test
suites can become extremely large—automated generation tools such as Randoop
and EvoSuite may produce thousands of test cases.  While these suites achieve
high code coverage and mutation scores, they often contain many redundant
tests, leading to long execution times and increased maintenance costs.  Test
suite minimisation aims to remove redundant tests while preserving the
effectiveness of the suite.  This thesis investigates how to perform test
suite minimisation directly on \emph{Stimulus–Response Matrices} (SRMs) within
the LASSO platform.

\section{Background and Motivation}
Regression testing verifies that software modifications do not break existing
functionality.  Software testing has become a major cost factor: early
studies estimate that it can consume nearly half of a software system's
development budget\cite{boehm1981}.  Large-scale test suites in modern
projects may contain thousands of automatically generated test cases and
take days or even weeks to execute\cite{rothermel1998}.  This escalating
complexity and cost demand effective strategies to keep test suites
manageable.  Test-suite minimisation offers a solution by removing
redundant tests to reduce execution time and maintenance overhead.  Yet
naïve reduction risks omitting tests that detect faults and thereby
undermines software quality\cite{yoo2012}.  Balancing reduction with
coverage and fault-detection retention is therefore crucial.

\section{LASSO Context}

The Large-Scale Software Observatorium (LASSO) is a research platform
developed at the University of Mannheim for large-scale software analysis.
It integrates static and dynamic analysis, automated test generation and
a domain-specific language for orchestrating analysis pipelines.  LASSO
records test executions in \emph{Stimulus–Response Matrices}.  An SRM is a
binary matrix where each row represents a test and each column represents a
coverage requirement (statement, branch or mutant); an entry of~1 indicates
that the test covers the requirement.  Although LASSO provides powerful
analysis capabilities, it currently lacks an integrated test-suite minimisation
component.  Adding minimisation would reduce redundant test execution and
improve pipeline throughput.

\section{Problem Statement}

Let~$T$ be a set of test cases and let~$R$ be the set of coverage
requirements.  Each test~$t\in T$ covers a subset $C(t)\subseteq R$.
Executing the entire suite~$T$ yields full coverage $C(T)=\bigcup_{t\in T}C(t)$.
The \emph{test suite minimisation} problem asks for the smallest subset
\(S^*\subseteq T\) such that $C(S^*)=C(T)$.  Solving this problem exactly is
NP-complete because it reduces to the set-cover problem.
Practical techniques therefore rely on heuristics and approximations that
balance reduction ratio, coverage retention and computational cost.

\section{Research Objectives and Questions}

The goal of this thesis is to design and implement an SRM-based test-suite
minimiser for LASSO.  To achieve this goal, we address the following
research questions:
\begin{enumerate}
  \item Which test-suite minimisation techniques are most widely studied and
        effective according to the literature?
  \item What criteria should be used to evaluate these techniques for
        integration into the SRM-based LASSO environment?
  \item Which technique best balances reduction ratio, coverage retention and
        scalability for use with LASSO?
\end{enumerate}

\section{Implementation Strategy}

Existing test-suite reduction tools (for example, the Java
Software Reduction framework) often support only a single programming
language or rely on language-specific program representations.
Because LASSO analyses software written in diverse languages, this thesis
adopts a language-independent strategy.  Rather than
integrating a pre-implemented program with limited language
support, we propose to implement the minimisation algorithm ourselves
directly on the SRM.  This approach enables reduction of test suites
generated for any language supported by LASSO and avoids dependencies
on external tools with fixed language bindings.

\section{Thesis Structure}

The remainder of this thesis is organised as follows:
\begin{description}
  \item[Chapter~2] introduces the fundamentals of code coverage, mutation
    testing, SRMs and formalises the test-suite minimisation problem.
  \item[Chapter~3] provides a comprehensive literature review of existing
    minimisation techniques, grouped into greedy heuristics, exact
    optimisation, search-based methods, data-driven approaches and empirical
    evaluations.
  \item[Chapter~4] defines evaluation criteria and compares candidate
    algorithms using these criteria.
  \item[Chapter~5] presents the rationale for selecting the Harrold–Gupta–Soffa
    algorithm with overlap-aware heuristics and outlines the proposed
    minimisation workflow.
  \item[Chapter~6] concludes the thesis and discusses future work.
\end{description}


